% !TeX document-id = {86da647f-a1fe-4e65-a2e3-4ae5a22335d2}
% !TEX program = xelatex
\documentclass[10pt,a4paper]{article}
\usepackage[top=13mm,bottom=12mm,left=15mm,right=15mm]{geometry}
\usepackage{fontspec}
\usepackage{xeCJK}
\usepackage{xcolor}
\usepackage{enumitem}
\usepackage{tabularx}
\usepackage{tcolorbox}
\usepackage{graphicx}
\usepackage{amsmath}
\usepackage{hyperref}
\tcbuselibrary{skins}

% Fonts available in the local environment.
\setmainfont[
  Path={arial/}, Extension=.ttf,
  UprightFont=ArialMdm, BoldFont=arialceb,
  ItalicFont=ArialMdmItl, BoldItalicFont=ArialCEBoldItalic
]{ArialMdm}
\setCJKmainfont[BoldFont={Source Han Sans TC Bold}]{Source Han Sans TC}
\setCJKmonofont{Source Han Sans TC}
\definecolor{navy}{HTML}{183348}
\definecolor{teal}{HTML}{167C80}
\definecolor{ink}{HTML}{263642}
\definecolor{muted}{HTML}{3F4F5B}
\definecolor{pale}{HTML}{F0F5F7}
\definecolor{rulecolor}{HTML}{CCDCE2}
\hypersetup{colorlinks=true,urlcolor=teal,pdftitle={楊子賢｜技術履歷},pdfauthor={楊子賢},pdfsubject={AI、機器學習與演算法}}
\pagestyle{empty}
\setlength{\parindent}{0pt}
\setlength{\parskip}{0pt}
\setlength{\emergencystretch}{2em}
\setlist[itemize]{leftmargin=1.15em,label=\textcolor{teal}{\textbullet},itemsep=2pt,topsep=3pt,parsep=0pt,partopsep=0pt}

\newcommand{\cvsection}[2]{%
  {\color{navy}\fontsize{12}{16}\selectfont\bfseries #1\hfill
  {\color{teal}\fontsize{7.2}{9}\selectfont #2}}\par
  \vspace{2pt}{\color{rulecolor}\hrule height 0.6pt}\vspace{5pt}}
\newcommand{\sideheading}[1]{\vspace{10pt}{\color{navy}\fontsize{10.5}{14}\selectfont\bfseries #1}\par\vspace{5pt}}
\newcommand{\award}[2]{\textbf{#1}\par{\color{muted}#2}\par\vspace{5pt}}
\newcommand{\projecttitle}[1]{{\color{navy}\fontsize{11.2}{15}\selectfont\bfseries #1}\par}
\newcommand{\tech}[1]{\vspace{2pt}{\color{teal}\fontsize{8.2}{11}\selectfont #1}\par}
% 大頭照：圓角裁切並加上細邊框
% \newcommand{\portraitwidth}{23mm}
% \newcommand{\portrait}[1]{%
%   \tcbincludegraphics[blankest,width=\portraitwidth,arc=1.6mm,
%     boxrule=0.8pt,colframe=rulecolor,
%     graphics options={width=\portraitwidth}]{#1}}

\begin{document}
\color{ink}

{\color{navy}\fontsize{31}{38}\selectfont\bfseries 楊子賢}\par
\vspace{3pt}
{\color{muted}\fontsize{10}{13}\selectfont ZIH-SIAN YANG}\par
\vspace{6pt}
{\color{muted}\fontsize{8.8}{11.8}\selectfont
國立臺灣海洋大學｜系所：資工系\par
\href{mailto:jihuty@gmail.com}{jihuty@gmail.com}｜\href{mailto:01257049@email.ntou.edu.tw}{01257049@email.ntou.edu.tw}\par
\href{https://github.com/YGJH}{github.com/YGJH}｜\href{https://www.kaggle.com/ygjh13}{kaggle.com/ygjh13}}

\vspace{6pt}
{\color{teal}\hrule height 2pt}
\vspace{8pt}

\fontsize{9.2}{12.3}\selectfont

% \begin{tcolorbox}[colback=pale,colframe=pale,boxrule=0pt,arc=2mm,
%   left=4mm,right=4mm,top=4mm,bottom=4mm,boxsep=0pt]
% {\color{navy}\fontsize{11.2}{14.8}\selectfont\bfseries 重點}\par
% \vspace{5pt}
% \textbf{\textcolor{teal}{Kaggle Expert}}\enspace
% {\color{muted}PTCG AI Battle｜銀牌｜237／6,807；Orbit Wars｜銅牌｜452／4,729}\par
% \vspace{4pt}
% \textbf{\textcolor{teal}{第一作者研究}}\enspace
% {\color{muted}EF1-Constrained Nash Social Welfare｜預印本；TAAI 2026 審查中}
% \end{tcolorbox}
\vspace{5pt}

\cvsection{競賽與檢定}{COMPETITIONS \& CERTIFICATIONS}
\textbf{\href{https://www.kaggle.com/ygjh13}{Kaggle Expert}}\par
\vspace{3pt}
\textbf{\href{https://www.kaggle.com/competitions/pokemon-tcg-ai-battle}{Kaggle PTCG AI Battle Challenge Simulation}}\enspace 237 / 6807 (3\%) \par
\vspace{3pt}
\textbf{\href{https://www.kaggle.com/competitions/orbit-wars}{Kaggle Orbit Wars}}\enspace 452 / 4729 (9\%) \par
\vspace{3pt}
\textbf{NCPC 2025}\enspace 決賽佳作\par
\vspace{3pt}
\textbf{TOPC 2025}\enspace （60／228）\par
\vspace{3pt}
\textbf{CPE 2025}\enspace 5 題｜76／2,515（前 3\%）\par
\vspace{3pt}
\textbf{TOEIC（多益）}\enspace 815 分｜藍色證書｜2026 年\par

\vspace{6pt}
\cvsection{研究成果}{RESEARCH HIGHLIGHT}
{\color{navy}\fontsize{9.8}{13}\selectfont\bfseries
EF1-Constrained Nash Social Welfare with Identical Additive Valuations: Complexity, Guarantees, and Experiments\par}
\vspace{4pt}
\textbf{第一作者}\enspace\textcolor{muted}{｜2026｜預印本；已投稿 TAAI 2026，審查中}\par
\vspace{3pt}
\textbf{主題：}本研究以不可分割物品分配為場景，同時追求 EF1 公平與 Nash Social Welfare（NSW），從理論界線一路做到可執行的學習式分配方法。\par
\begin{itemize}
\item \textbf{回答公平的代價：}證明問題為 strongly NP-complete；即使要求 EF1，仍保證至少保留最優 NSW 的 \textbf{69.2\%}，而當單一物品相對夠小時，效率更會\textbf{趨近最優}。
\item \textbf{讓公平能逐步維持：}證明在 identical valuations 下，每次將新物品分給目前效用最低者，即可在整個分配過程中持續滿足 EF1。
\item \textbf{把定理變成模型保證：}提出 PriorityNet，以 PPO 學習分配策略，並用 action mask 排除不公平動作；\textbf{公平由機制保證，模型專注提升福利}。
\item \textbf{3,000 筆實例驗證：}offline／online mean normalized NSW 達 \textbf{0.9911／0.9701}；相較 LPT／greedy baseline，win-minus-loss 為 \textbf{+27.10\%／+17.87\%}，online large regime 達 \textbf{+35.40\%}。
\end{itemize}
\href{https://arxiv.org/abs/2609.03846}{\textbf{arXiv:2609.03846}}

\vspace{5pt}
\cvsection{專案}{SELECTED PROJECTS}

\projecttitle{Pokémon TCG 決策代理人}
{\color{muted}Transformer 模仿學習系統｜\href{https://www.kaggle.com/competitions/pokemon-tcg-ai-battle}{Kaggle PTCG AI Battle} 銀牌}\par
\tech{PyTorch · Transformer · Pointer Network · Behavior Cloning · W\&B}
\begin{itemize}
\item \textbf{Kaggle 銀牌（237／6,807）：}以端到端決策代理人直接銜接官方 C++ 遊戲引擎，完成競賽部署。
\item \textbf{準確率成果：}兩個牌組的 specialist 相對基準準確率增益達 \textbf{+20.5／+49.4 個百分點}；以 Transformer 編碼 46 個狀態 token，並為動態合法選項評分。
\item \textbf{可擴展訓練：}以數百萬個決策點訓練，整合 memory-mapped shard、GPU 特徵處理、bfloat16、EMA 與 Muon optimizer。
\end{itemize}

\vspace{4pt}
\projecttitle{Orbit Wars 策略代理人}
{\color{muted}動態幾何規劃與多波次策略｜\href{https://www.kaggle.com/competitions/orbit-wars}{Kaggle Orbit Wars} 銅牌}\par
\tech{Rust · Python · 啟發式搜尋 · 動態幾何規劃 · 貪婪最佳化}
\begin{itemize}
\item \textbf{Kaggle 銅牌（452／4,729）：}設計自主策略代理人，在行星持續移動與敵方可能增援的環境中規劃攻防。
\item \textbf{未來狀態規劃：}預測軌道位置並求解 ETA 攔截；以抵達時的守軍、敵方增援風險與行星產能評估攻擊價值。
\item \textbf{即時決策與部署：}以 ROI 導向的貪婪法選擇多波次行動、保留前線防禦；採 Rust 擴充正式推論。
\end{itemize}


\vspace{5pt}
\cvsection{學歷}{EDUCATION}
\textbf{國立臺灣海洋大學}\enspace
{\color{muted}資訊工程學系}\par

\vspace{5pt}
\cvsection{技能}{SKILLS}
\textbf{ML／研究：}\enspace PyTorch、Transformer、PPO、W\&B、模仿學習\par
\vspace{3pt}
\textbf{興趣：}\enspace 啟發式搜索、Python、Rust、C/C++、Autodesk 模型建模\par

\end{document}
